\documentclass[ngerman,a4paper,11pt,twoside,cleardoubleempty,halfparskip]{scrreprt}

% Erforderliche Pakete
\usepackage[utf8]{inputenc}
\usepackage{babel}
\usepackage{graphicx}
\usepackage{fancyhdr}
\usepackage{multirow}
\usepackage{ifthen}
\usepackage{calc}
\usepackage{tabularx}
\usepackage{listings}
\usepackage[autokw]{svn-multi}
\usepackage[toc]{glossaries}
\usepackage[implicit=false]{hyperref} % implizit stört irgendwie beim Lesen

% weiteren, nützliche Pakete
%\usepackage{subfigure} 
%\usepackage{varioref}
%\usepackage{framed}
%\usepackage{longtable}
%\usepackage{floatflt}
%\usepackage{amsmath}
%\usepackage{url}
% ...

% Pakete für Verzeichnisse. Diese sind hier aufgeführt, da mindestens das
% 'glossaries' Paket nach z.B. hyperref geladen werden muss.
\usepackage{bibgerm}

\svnid{$Id: spec.tex 107 2012-05-12 16:23:28Z dgens001 $}

% Glossar(e) laden und erstellen
\makeglossaries
\loadglsentries{meta/glossar} % Name der Glossardatei (ohne .tex)

% Befehlsdefinitionen
\newcommand{\clearemptydoublepage}{\newpage{\pagestyle{empty}\cleardoublepage}}

% Globale Formatierungseinstellungen
%\renewcommand*\cleardoublepage{\clearpage\if@twoside \ifodd\c@page
%  \hbox{}\newpage\if@twocolumn\hbox{}\newpage\fi\fi\fi}
\renewcommand{\encodingdefault}{OT1}
\renewcommand{\familydefault}{cmss} % Schriftfamilie auf Sans Serif
\renewcommand{\glsdisplayfirst}[4]{\textit{#1#4}}
\setcapindent{1em}

% Anpassung der Ränder und Breitenverhältnisse
% Der letzte Wert in oddsidemargin bestimmt die Bindekorrektur
\setlength{\oddsidemargin}{2cm - 1in + 0.5cm}
\setlength{\textwidth}{\paperwidth - (1in + \hoffset) - \oddsidemargin - 4cm}
\setlength{\evensidemargin}{\paperwidth - (1in + \hoffset)*2 - \oddsidemargin - \textwidth}
\setlength{\marginparwidth}{4cm - \marginparsep{} - 1cm}

% Wenn größerer Zeilenabstand gewünscht ist, je nach Schriftgröße:
% Abstand        10pt    11pt    12pt
% -----------------------------------
% anderthalb     1.25    1.21    1.24
% doppelt        1.67    1.62    1.66
%
%\renewcommand{\baselinestretch}{1.21}

% Kopf- und Fußzeilen einrichten
\pagestyle{plain}
\setlength{\headwidth}{\textwidth}
\addtolength{\headwidth}{\marginparsep}
\addtolength{\headwidth}{\marginparwidth}
\renewcommand{\chaptermark}[1]{\markboth{#1}{}}
\renewcommand{\sectionmark}[1]{\markright{\thesection\ #1}}
\lhead[\fancyplain{}{\bfseries\thepage}]%
	{\fancyplain{}{\bfseries\rightmark}}
\rhead[\fancyplain{}{\bfseries\leftmark}]%
	{\fancyplain{}{\bfseries\thepage}}
\cfoot{}

\begin{document}
\pagenumbering{arabic}
\svnid{$Id: titel.tex 180 2012-05-24 15:26:04Z dgens001 $}
\begin{titlepage}
	\begin{center}
		% Kopf der Seite
		\includegraphics{meta/hsrm-logo} \\[0.7cm]
		{Fachbereich Design Informatik Medien}
		
		\vfill

		% Titel
		%\rule{\textwidth}{1pt}\\[0.5cm]
		{\huge \bfseries Designspezifikation Campusadventure}\\[0.1cm]
                im Fach Softwaretechnik im SS2012
		%\rule{\textwidth}{1pt}
		
		\vfill
		% Mitte der Seite
		\begin{tabular}{lr}
			Vorgelegt von & \line(1,0){100} \\
						  & Dominik Schuhmann \\ 
						  & \\
						  & \line(1,0){100} \\
                          & Tino Landmann \\ 
						  & \\
						  & \line(1,0){100} \\
                          & Simon Hardt \\
						  & \\
						  & \line(1,0){100} \\
                          & David Gens \\
                        \\
						\\
                        \\
						\\
			bei           & Prof. Dr. Wolfgang Weitz \\
                        \\
						am			  & \today \\
						\\
                        in der Version       &	\svnrev. \\
		\end{tabular}
		\vfill
		
		% Fuß der Seite
		
	\end{center}
\end{titlepage}
 % Titelseite (wird automatisch mit Werten von oben gefüllt)
\clearemptydoublepage % Sorgt dafür, dass die nächste Seite rechts beginnt

  \begin{longtable}{l l l p{80px} p{200px}}
	  \textbf{Rev} & \textbf{Autor}  &\textbf{Datum} & \textbf{Beschreibung} & \textbf{Dateien}\\
	  \hline
	              &      \           &       \       &            \          &        \\
	  181 &
		  tland001
		  &
		  24-05-2012
		  &
		  lautfzeitdiagram bla2
		  & \strut
		  /design/kapitel/laufzeitsicht.tex \par  \strut \\
	  180 &
		  dgens001
		  &
		  24-05-2012
		  &
		  svn tag mal gesetzt.. wäre clever gewesen
		  & \strut
		  /design/kapitel/bausteinsicht.tex \par /design/meta/svnlog.tex \par /design/kapitel/musterundkonzepte.tex \par /design/meta/glossar.tex \par /design/design.pdf \par /design/meta/titel.tex \par /design/kapitel/laufzeitsicht.tex \par /design/kapitel/einfuehrung.tex \par /design/design.tex \par /design/kapitel/entwurfsentscheidungen.tex \par  \strut \\
	  179 &
		  dgens001
		  &
		  24-05-2012
		  &
		  laufzeitdiagramme umbenannt und beispielhaft eingefuegt
		  & \strut
		  /design/kapitel/laufzeitsicht/Spiel Laden.png \par /design/kapitel/laufzeitsicht/spielLaden.png \par /design/kapitel/bausteinsicht.tex \par /design/meta/svnlog.tex \par /design/kapitel/laufzeitsicht/Spieler\_Steuerung.png \par /design/kapitel/musterundkonzepte.tex \par /design/design.pdf \par /design/kapitel/laufzeitsicht/Neues Spiel.png \par /design/kapitel/laufzeitsicht/neuesSpiel.png \par /design/kapitel/laufzeitsicht.tex \par /design/bewegen\_seq.vpp \par  \strut \\
	  
  \end{longtable}
   % die letzten paar SVN-Logs als TEX-Tabelle formatiert
                    % (siehe meta/svnlog-gen.sh und meta/svnlog.xslt)
\clearemptydoublepage
\tableofcontents % Inhaltsverzeichnis generieren
\clearemptydoublepage

% --------------------------------- Hauptteil ----------------------------------
% Im diesen Teil werden die einzelnen Kapitel eingefügt, die sinnvollerweise
% im Verzeichnis 'kapitel' abgelegt werden.
\svnid{$Id: einleitung.tex 102 2012-04-30 07:46:35Z shard001 $}
\chapter{Einleitung}
In folgendem Dokument werden die Grundlagen und Anforderungen des Softwaretechnikprojektes definiert und 
spezifiziert. 

Abschnitt 2 bietet eine Übersicht über die grundlegenden Anforderungen an das Projekt, sowie die technischen 
Anforderungen und unsere Umsetzungsideen. Der Abschnitt 3 befasst sich mit den allgemeinen Abläufen und 
Funktionen der Anwendung, die sich in Menü- und Ingamemodus aufteilen. Einen Überblick über das Domänenmodell 
zeigt Abschnitt 4 anhand einer genauen Beschreibung der einzelnen Elemente mit Hilfe eines UML-Diagrammes. 
Die Benutzerschnittstelle zum Spieler wird in Abschnitt 5 mittels einiger Zeichnungen beschrieben. Ebenso 
ist das Campusdateiformat hier definiert.


\svnid{$Id: grundlagen.tex 108 2012-05-12 18:44:01Z dgens001 $}
\chapter{Projektgrundlagen}

\textit{Dieses Kapitel bietet eine kurze Übersicht über definierte Ziele und Rahmenbedingungen. Es bietet
einen ersten Architekturüberblick und stellt grobe Systembestandteile mit ihren Zuständigkeiten dar.}

\section{Zielsetzung}
Ziel des Projekts ist eine selbst entwickelte Spieleapplikation, die ohne
aufwändige Grafikdarstellung oder KI einen amüsanten Zeitvertreib bieten
kann. Nicht ganz nebenbei soll die Applikation Ergebnis eines größeren
Projekts sein, an dem sich softwaretechnische Methoden, kollaborative Arbeit
und der nahezu vollständige Entwicklungszyklus üben und testen lassen.
Fachliche, technische und andere qualitative Ziele sollen in diesem Dokument
zusammengefasst werden.

\section{Inhaltlicher Überblick}
Die Autoren dieser Spezifikation haben sich darauf verständigt, ein Spiel nach der Machart einiger
bereits bekannter Roleplaying-Adventure zu entwickeln. Das Spiel könnte also (je nach Storyline) endlos 
spielbar sein, oder ein festes Ende haben, wobei das hier spezifierte Spiel einen definiertes Spielende
haben wird.
Um einen gewissen Anreiz für den \gls{Spieler} zu schaffen, sollte die Storyline 
diesen relativ frei durch den \gls{Campus} führen. Dabei wird es in späteren Erweiterungen
möglich sein, den eigenen \gls{Charakter} zu leveln. Inhaltliches Ziel soll das Finden
eines bestimmten \gls{Schein}s sein.
Der Spieler befindet sich immer auf genau einem \gls{Feld}. Er kann nur mit \glspl{Item} oder \gls{npcs}, die
sich auf dem gleichen \gls{Feld} befinden interagieren. Zudem hat er die Möglichkeit zu einem direkt 
benachbarten \gls{Feld} zu wechseln. Mehrere Felder bilden einen \gls{Raum}. \glspl{Raum} sind in sich 
geschlossene Areale. Räume haben einen oder mehrere \glspl{Eingang} und müssen dem \gls{Spieler} nicht 
zugänglich sein. Alle Räume zusammen bilden den \gls{Campus}. 
Von Zeit zu Zeit sollte der \gls{Spieler} einen oder mehrere Gegenstände (sog. \glspl{Item}) erhalten, die er
in seinen \gls{Inventar} aufnehmen kann. Manche \glspl{Item} sind wertvoll (hilfreich für den Verlauf des 
Spiels), andere nicht. Einige \glspl{Item} sollten sich kombinieren lassen, oder in anderer Form mit 
dem \gls{Spieler} oder der \gls{Welt} interagieren können. Der \gls{Spieler} sollte ein \gls{Savegame} 
anlegen und laden können. Das \gls{Savegame} beinhaltet die Campuskonfiguration und alle Elemente, die sich im
Spiel befinden. Es hat ein Datum und eine Versionsnummer (für den Fall einer Erweiterung).

\section{Technische Anforderungen}
Die Applikation soll unter Java Swing entwickelt werden. Sie sollte auf 
aktuellen Desktop-PCs ohne Leistungsprobleme spielbar sein (10MB freier 
Festplattenspeicher, 1GB RAM, 2GHz DualCore, 128MB Grafikspeicher, 
Bildschirmauflösung 1280x1024). Die Applikation sollte im Fenstermodus mit einer Fenstergröße von 1024x768
laufen und nicht vergrößer- oder verkleinerbar sein (entspricht Seitenverhältnis 4:3).
Der Fokus soll aber nicht auf möglichst realistischer Echtzeitgrafik, sondern auf dem Spielerlebnis liegen.
Darüber hinaus sollte der Spielplan anhand einer Spielplandatei konfigurierbar sein. 
Dort sollte man zunächst Wände und Spielstart, später vielleicht auch \acrlong{npc} und \glspl{Item} 
positionieren können.

\section{Architektur}
Die Anwendungslogik soll derart beschaffen sein, dass der Spieler anhand einfacher Zugriffe (3-4 Klicks)
zu den gesuchten Funktionen findet. Ziel ist nicht, dem Spieler über möglichst viele Einstellungsmöglichkeiten
verfügen zu lassen. Es gibt ein Itemsystem, ein Raumsystem und ein Dialogsystem. In eventuellen Erweiterungen
können zusätzliche Systeme implementiert werden. Die Architektur sollte später in Darstellung, Logik und 
Anwendungsdaten systematisch unterteilt werden. Generell lässt sich das Spiel inhaltlich in zwei Modi 
unterteilen, Ingame- und \gls{Menuemodus}. Der Menümodus sollte der Standardmodus sein, in dem die Applikation
startet und von dem aus man in den Ingamemodus gelangt. Der Einfachheit halber sind die Abbruchsübergänge in
Abbildung \ref{fig:zustaende_entwurf} nicht berücksichtigt.

\begin{figure}[h]
	\begin{center} 
		\includegraphics[width=150mm, height=60mm]
                  {kapitel/grundlagen/architektur.png}
	\end{center}
	\caption{grobe Struktur der Systemzustände}
	\label{fig:zustaende_entwurf}
\end{figure}

\section{Erweiterbarkeit}
Der \gls{Charakter} kann in späteren Erweiterungen z.B. verschiedene Fähigkeiten oder
Eigenschaften haben. Diese könnten dann im Laufe des Spiels weiterentwickelt werden 
(z.B. anhand sog. CreditPoints). CreditPoints würde man durch das erfolgreiche Sammeln
der \glspl{Schein} erhalten, mehrere \glspl{Schein} könnten dann ein Semester ergeben.
Es könnte dann auch einen Abschluss in Form eines höchsten Semesters (z.B. Master of Sciene) geben.
Wahlweise könnte der \gls{Spieler} dann trotz Erreichen dieses Semesters weiterhin 
durch die \gls{Welt} navigieren und zufällig auftauchende \glspl{Schein} (z.B. durch Lösen von
Minispielen oder Rätseln) finden. Es könnte dann auch die Möglichkeit geben, unabhängig von den 
\glspl{Savegame} Charakterprofile zu speichern, zu editieren und zu laden. Daher soll es möglich sein, 
die Gegenstände, Objekte und Charaktere im Spiel von Version zu Version zu ergänzen und zu erweitern. 
So könnten in einer zukünftigen Erweiterung der NPCs feste Routen vorgesehen werden, auf denen diese sich
wärend dem Spiel bewegen.

\svnid{$Id: funktionen.tex 108 2012-05-12 18:44:01Z dgens001 $}
\chapter{Abläufe und Funktionen}

\textit{Dieses Kapitel bietet einen Einblick in das Zusammenspiel der Prozesse und Abläufe des Spiels.
Anwendungsfälle der Interaktionen mit dem System, dabei wichtige Anwendungsfunktionen und die jeweiligen 
Akteure sollen einen klaren Umriss des später implementierten Spiels geben.}

\section{Menümodus}
Nach dem Start der Anwendung findet sich der \gls{Spieler} im Menümodus. Von dort aus kann er die Funktionen
des Hauptmenüs erreichen und in andere Menüsichten wechseln. Unter Optionen könnte eine Reihe von 
Einstellungen getroffen werden, zumindest soll hier die Ingamesteuerung konfigurierbar sein. Später könnten
auch zusätzliche Sound- oder Videoeinstellungen hier vorgenommen werden. Die Funktion \textit{Spiel beenden}
beendet die Applikation.

\begin{figure}[htb]
	\begin{center} 
		\includegraphics[width=100mm, height=60mm]
                  {kapitel/funktionen/Hauptmenue.png}
	\end{center}
	\caption{Übersicht der Hauptmenü-Funktionen}
	\label{fig:hauptmenue_funktionen}
\end{figure}

Neben dem Menümodus mit seinen verschiedenen Sichten gibt es den Ingamemodus, in dem das eigentliche Spiel
abläuft. Die Funktionen \textit{Spiel laden} und \textit{Neues Spiel starten} wechseln vom Menü- in den
Ingamemodus. Startet man ein neues Spiel, muss zunächst ein \gls{Campus} geladen werden. Hierzu wählt man 
aus einer zuvor erstellten Liste eine Campusdatei aus. Die \gls{Welt} wird anhand dieser Datei automatisch 
erstellt. Bestätigt der Spieler seine Campuswahl, wechselt das System automatisch in den Ingamemodus und
startet auf dem gewählten Campus.

\paragraph{UseCase: Neues Spiel starten}
\begin{center}
	\begin{tabular}{|r l|}
	  \hline
	  \textbf{Titel}: & Neues Spiel starten \\
	  \textbf{Akteur}: & \gls{Spieler} \\
	  \textbf{Fachlicher Auslöser}: & \gls{Spieler} hat \textit{Neues Spiel starten} gewählt \\
	  \textbf{Vorbedingung}: & Das System hat Zugriff auf mindestens eine Campusdatei \\
	  \hline
	  \textbf{Standardablauf}:
		& 1. Spieler: Wählt einen Campus aus \\
		& 2. Spieler: Bestätigt Campuswahl \\
		& 3. System: Liest die Campusdatei ein \\
	  \textbf{Abbruch}:
		& 1a. Spieler: Wählt \textit{Zurück} \\
		& 2a. System: Wechselt ins Hauptmenü \\
	  \hline
	  \textbf{Nachbedingung}: & System geht in Zustand \gls{Ingame} oder Hauptmenü \\
	  \textbf{NF-Anf}.: & Ladezeit $<$ 30 Sekunden \\
	  \textbf{Parameter}: & Campusdatei \\
	  \textbf{Nutzungshäufigkeit}: & bei jedem neuen Spiel \\
	  \hline
	\end{tabular}
\end{center}
\newpage
\section{Ingamemodus}
Der Ingamemodus ist der wichtigste Zustand der Applikation. Hier bieten sich dem Spieler alle Möglichkeiten,
er kann den Campus erkunden, Gegenstände aufnehmen oder ablegen, mit NPCs reden, neue Räume betreten usw.
Im Wesentlichen lassen sich die Funktionen wie in Abb. \ref{fig:ingame_funktionen} dargestellt zusammenfassen.

\begin{figure}[htb]
	\begin{center} 
		\includegraphics[width=100mm, height=50mm]
                  {kapitel/funktionen/InGame.png}
	\end{center}
	\caption{Übersicht der Ingame-Funktionen}
	\label{fig:ingame_funktionen}
\end{figure}

\subsection{Charakterbewegung}
Der \gls{Spieler} kann sich grundsätzlich über entsprechende Tasten von Feld
zu Feld bewegen oder seine Sicht innerhalb eines Feldes steuern. Dabei behält der
\gls {Spieler} bei seitlichen Bewegungen über Felder stets seine Blickrichtung bei (sogenanntes Strafing).
Die Blickrichtung kann er ändern, indem er sich auf einem Feld nach Links oder Rechts dreht.
Grundsätzlich gibt es nur vier Blickrichtungen auf jedem Feld, die mit Nord, Süd, Ost und
West bezeichnet werden. Bei Bewegung nach vorn oder hinten behält der Spieler ebenfalls seine Blickrichtung
bei (die Ansicht wechselt also nicht um $180 ^{\circ}$ wenn der Spieler rückwärts gehen möchte). 

\subsubsection{UseCase: Charakterbewegung}
\begin{center}
	\begin{tabular}{|r l|}
	  \hline
	  \textbf{Titel}: & Charakterbewegung\\
	  \textbf{Akteur}: & Spieler \\
	  \textbf{Fachlicher Auslöser}: & Spieler wählt Bewegungsrichtung \\
	  \textbf{Vorbedingungen}: & Spieler ist Ingame \\
	  \hline
	  \textbf{Laufen}:
		& 1. Spieler: Möchte nach vorn oder hinten laufen \\
		& 2. System: Versetzt den Spieler auf das Feld seiner Wahl \\
		& 3. System: Wechselt die Ansicht zum neuen Feld \\
	  \textbf{Strafen}:
		& 1. Spieler: Möchte nach rechts oder links strafen \\
		& 2. System: Versetzt den Spieler auf das Feld seiner Wahl \\
		& 3. System: Wechselt die Ansicht zum neuen Feld \\
	  \textbf{Drehen}:
		& 1. Spieler: Möchte Blickrichtung nach links oder rechts wechseln \\
		& 2. System: Dreht die Ansicht auf dem aktuellen Feld \\
		& um $90\,^{\circ}$ in die gewählte Richtung\\
	 \textbf{Fehlerfall Bewegen}:
		& 2.a System: Feld ist dem Spieler nicht zugänglich (z.B. Wand) \\
		& 3.a System: Mitteilung an den Spieler, Ansicht wechselt nicht \\
	  \hline
	  \textbf{Nachbedingung}: & Ansicht gewechselt (im Fehlerfall nicht) \\
	  \textbf{NF-Anf}.: & Reaktionszeit System $<$ 2 Sekunden \\
	  \textbf{Parameter}.: & Feld, Blickrichtung, Ansicht \\
	  \textbf{Nutzungshäufigkeit}: & Ingame sehr häufig (ca. 5 von 10 Interaktionen) \\
	  \hline
	\end{tabular}
\end{center}

\subsection{Interaktion mit Items}
\glspl{Item}, die sich im Besitz des \gls{Spieler}s befinden, können verwendet oder miteinander kombiniert 
werden. Dazu muss der \gls{Spieler} ein \gls{Item} aus dem Inventar in die Hand nehmen und über die 
Benutzeninteraktion
mit einem Gegenstand aus dem Inventar kombinieren oder auf Objekte oder NPCs auf dem Spielfeld 
anwenden. Häufig wird, wenn ein \gls{Item} benutzt wird, dieses verbraucht werden und aus dem \gls{Inventar}
verschwinden. Es gibt jedoch auch viele Gegenstände, welche mehrmals verwendet werden können, die nicht nach
ihrer Benutzung verschwinden. Um den begrenzten \gls{Inventar}platz zu verwalten, können \glspl{Item} auch
auf dem 
aktuellen \gls{Feld} abgelegt werden. Dazu wird es in die Hand genommen und dann über die Drop-Funktion 
abgelegt. Das \gls{Item} ist dann auf dem Feld wieder sichtbar und kann aufgenommen werden. Dazu wählt der 
\gls{Spieler} in der Infosicht den gewünschten Gegenstand aus, diesen hat er dann wieder in der Hand 
und kann entscheiden ob er ihn direkt benutzen, im Inventar ablegen oder wegwerfen möchte.

\paragraph{Szenario: Item finden}
Der \gls{Spieler} bewegt sich frei über den Campus. Dabei liegen auf dem Feld vor ihm drei \glspl{Item}. Der 
Spieler möchte nun das \gls{Item} aufnehmen. Es öffnet sich eine Liste, in der alle \glspl{Item} aufgelistet 
werden. Der \gls{Spieler} browst nun durch die Liste und wählt einen Schlüssel aus. Die anderen \glspl{Item}
braucht er nicht. Jedoch bekommt er, bevor er den Gegenstand aufgenommen hat, eine Meldung das sein Inventar
bereits voll ist und der Schlüssel daher nicht aufgenommen werden kann. Aus diesem Grund wirft er das Sandwich
aus seinem Inventar. Das Sandwich erscheint nun auf dem Feld und ein Slot im Inventar ist frei geworden. Nun 
kann er den Schlüssel aufnehmen. Der Schlüssel ist dabei vom Boden verschwunden und stattdessen nun in seinem 
Inventar.

\paragraph{UseCase: Item finden}
\begin{center}
	\begin{tabular}{|r l|}
	  \hline
	  \textbf{Titel}: & Item Finden \\
	  \textbf{Akteur}: & Spieler \\
	  \textbf{Fachlicher Auslöser}: & Spieler schaut sich um \\
	  \textbf{Vorbedingungen}: & Spieler ist Ingame auf einem Feld mit Items\\
	  \hline
	  \textbf{Standardablauf}:
		& 1. Spieler: Schaut sich um und bekommt alle Items auf dem Feld,\\
	    & in Form einer Liste (in der Infosicht) angezeigt \\
		& 2. Spieler: Wählt ein Item aus dieser Liste\\
		& und nimmt dieses in sein Inventar auf \\
		& 3. System:  Übergibt Item von Feld an Spieler \\ 
	 \textbf{Abbruch}:
                & 2.a Spieler: Möchte kein Item aufnehmen und bricht ab\\
                & 3.a System: Wechselt die Anzeige von Infosicht zu Inventarsicht\\
	  \hline
	  \textbf{Nachbedingung}: & Spieler erhält Item \\
	  \textbf{NF-Anf}.: & Reaktionszeit System $<$ 2 Sekunden \\
	  \textbf{Parameter}: & Inventar / Item \\
	  \textbf{Nutzungshäufigkeit}: & \gls{Ingame} häufig (ca. 3 von 10 Interaktionen) \\
	  \hline
	\end{tabular}
\end{center}

\paragraph{UseCase: Inventarinteraktion}
\begin{center}
	\begin{tabular}{|r l|}
	  \hline
	  \textbf{Titel}: & Inventarinteraktion \\
	  \textbf{Akteur}: & \gls{Spieler} \\
	  \textbf{Fachlicher Auslöser}: & \gls{Spieler} möchte einen seiner Gegenstände verwenden \\
	  \textbf{Vorbedingungen}: & Im \gls{Inventar} befinden sich verwendbare Gegenstände \\
	  \hline
	  \textbf{Ablauf}:
		& 1. Spieler: Wählt sein Inventar an \\
		& 2. System: Inventarsicht wird aufgerufen \\
		& 3. Spieler: Wählt einen Gegenstand \\
		& 4. System: Wechselt zur Infosicht \\

	  \hline
	  \textbf{Nachbedingung}: & Bei leerem Handslot wird dieser mit Item gefüllt und \\
	  	& Gegenstand wird aus dem Inventar entfernt, bei besetztem \\
		& Handslot wird das Item im Handslot ersetzt und \\
		& wandert in das Inventar \\
	  \textbf{NF-Anf}.: & Reaktionszeit System $<$ 2 Sekunden \\
	  \textbf{Parameter}: & ausgewählter Gegenstand \\
	  \textbf{Nutzungshäufigkeit}: & beliebig oft, sofern sich Gegenstände im Inventar befinden \\
	  \hline
	\end{tabular}
\end{center}

\subsection{Interaktion mit NPCs}
Befindet sich ein \gls{npc} auf dem Feld des Spielers, kann der Spieler über die Interaktionssicht mit 
diesem in einen
Dialog treten. Ein Dialog wird einen kurzen Einleitungstext haben. Im Verlauf des Dialogs kann der NPC
dem \gls{Spieler} Handlungsanweisungen geben, ihm Fragen stellen oder Items übergeben. Dem Spieler werden
dabei mehrere Frage- oder Antwortmöglichkeiten zur Verfügung gestellt. Gewisse Fragen oder Antworten 
treiben den Dialog voran, während andere ihn zum Ende bringen.
Bestimmte Dialoge können \glspl{Item} beinhalten. Wird ein solcher Dialog erreicht, vergleicht das 
System das (Ist)-\gls{Item}, welches der \gls{Spieler} dem NPC anbietet mit dem (Soll)-\gls{Item}, dass
der NPC verlangt um eine bestimmte Aufgabe als erfüllt anzuerkennen. Es kann jedoch auch
nach dem Erfüllen einer Aufgabe bzw. damit eine Aufgabe erfüllt wird, der \gls{Spieler} ein bestimmtes 
\gls{Item} vom NPC erhalten.

\paragraph{UseCase: Interaktion mit NPC}
\begin{center}
	\begin{tabular}{|r l|}
	  \hline
	  \textbf{Titel}: & Interaktion mit NPC \\
	  \textbf{Akteur}: & \gls{Spieler} und \gls{npc} \\
	  \textbf{Vorbedingungen}: & \gls{Spieler} steht auf gleichem \gls{Feld} wie der \gls{npc} \\
	  \hline
	  \textbf{Ablauf}:
		& 1. Spieler: \gls{Spieler} spricht \gls{npc} an \\
		& 2. System:  \gls{Dialog} des \gls{npc} wird aufgerufen \\
		& 3. Spieler: \gls{Spieler} wählt eine \gls{Erwiderung} aus \\
		& 4. System:  Wird fortgeführt bis \gls{Dialog} abgeschlossen \\
	  \textbf{Ablauf (mit Item)}:
		& 1. Spieler: \gls{Spieler} wendet Gegenstand auf \gls{npc} an \\
		& 2. System:  \gls{npc} reagiert ggf., je nach Art des Gegenstandes \\
		& 3. Spieler: \gls{Spieler} wählt eine \gls{Erwiderung} aus, sofern möglich \\
		& 4. System:  Wird fortgeführt bis \gls{Dialog} abgeschlossen oder beendet \\
	  \hline
	  \textbf{Nachbedingung}: & \gls{Dialog} mit \gls{npc} wird an aktuellem Stand angehalten \\
	  \textbf{NF-Anf}.: & Systemseits 1-2 Sekunden Reaktionszeit \\
	  \textbf{Parameter}: & \gls{npc}, evtl ausgewählter Gegenstand \\
	  \textbf{Nutzungshäufigkeit}: & Ingame häufig (3 von 10 Interaktionen)\\
	  \hline
	\end{tabular}
\end{center}

\subsection{Pausieren}
Pausiert der \gls{Spieler} das laufende Spiel, so wechselt er kurzzeitig in den Menümodus mit
der Pausesicht, wo er z.B. seinen Spielstand speichern oder einen alten laden kann. Das Spiel 
selbst bleibt, während er im Pausemenü ist, erhalten und kann jeder Zeit fortgesetzt werden. 
Über den Pausemodus kann zudem das Spiel abgebrochen werden. Dabei gelangt der \gls{Spieler} 
zurück ins Hauptmenü. Das Pausemenü ist gleichzeitig die einzige
Möglichkeit für den Spieler, die Applikation aus dem Ingamemodus heraus regulär zu beenden.
\begin{figure}[h]
	\begin{center} 
		\includegraphics[width=90mm, height=58mm]
                  {kapitel/funktionen/Pausemenue.png}
	\end{center}
	\caption{Übersicht der Pausemenü-Funktionen}
	\label{fig:pausemenue_funktionen}
\end{figure}

\newpage

\paragraph{UseCase: Spiel pausieren}
\begin{center}
	\begin{tabular}{|r l|}
	  \hline
	  \textbf{Titel}: & Spiel pausieren\\
	  \textbf{Akteur}: & Spieler \\
	  \textbf{Fachlicher Auslöser}: & Spieler möchte Spiel unterbrechen \\
	  \textbf{Vorbedingungen}: & Spieler ist \textit{Ingame} \\
	  \hline
	  \textbf{Standardablauf}:
		& 1. Spieler: wählt Pausenmenü \\
		& 2. System: wechselt in den \textit{Menümodus} \\
		& 3. Spieler: wählt eine Pausefunktion aus \\
	  \textbf{Abbruch}:
                & 1.a Spieler: Spieler wählt \textit{Spiel fortsetzen} \\
		& 2.a System:  wechselt in den \textit{Ingamemodus} \\
	  \hline
	  \textbf{Nachbedingung}: & Spiel pausiert, System im Menümodus\\
	  \textbf{NF-Anf}.: & Reaktionszeit System $<$ 1 Sekunde \\
	  \textbf{Nutzungshäufigkeit}.: &  min. einmal pro Spiel \\
	  \hline
	\end{tabular}
\end{center}

\subsection{Speichern eines Spiels}
Das Spiel kann ausschließlich gespeichert werden, wenn der \gls{Spieler} sich im
Ingamemodus befindet. Dazu wählt er die entsprechende Taste zur Spielunterbrechung.
Darüber gelangt er dann in das Pausemenü. An dieser Stelle gibt es verschieden Auswahloptionen 
u.a. die das Spiel zu speichern. Geht der \gls{Spieler} in das Spiel speichern
Menü, kann er an dieser Stelle einen vorhandenen Spielstand aus einer Liste 
überschreiben oder einen neuen anlegen und so sein derzeitiges Spiel speichern. Ihm 
steht es jedoch frei auch einen Spielstand zu löschen oder ohne zu speichern das Menü 
zu verlassen.

\subsection{Laden eines Spiels}
Spiele können im CampusAdventure auf zwei Arten geladen werden. Zum einen direkt nach dem Starten
der Spieldatei. Der \gls{Spieler} erreicht dabei zunächst das Hauptmenü, in dem er neben anderen
Optionen, die des Spiel ladens hat. Das Spiel laden Menü selbst, ist vom Aufbau her gleich dem
Speichermenü. Es kann ein Spielstand aus der Liste gewählt werden und dieser wird entweder gelöscht,
geladen oder das Menü wird ohne jegliche Aktion wieder verlassen, was einen Abbruch darstellt. Die 
zweite Variante ist, dass während einer Spielunterbrechung ein anderer Spielstand geladen werden kann.
Hierzu befindet sich der \gls{Spieler},auf Grund der Spielunterbrechung, im Pausemenü. In diesem ist
dann das Spiel laden Menü integriert und nutzbar.
\newpage
\
\svnid{$Id: daten.tex 108 2012-05-12 18:44:01Z dgens001 $}
\chapter{Daten- und Domänenmodell}
\textit{In diesem Kapitel bemühen wir uns um eine formale Darstellung der Gegenstandswelt der Anwendung. Es 
werden die Objekthierarchie und die wichtigsten Konzepte, die in der \gls{Welt} auftauchen, beschrieben.}

\section{Überblick}
Im \gls{Spiel} gibt es \glspl{GameObject} und einen \gls{Campus} (später eventuell mehrere). 
Der \gls{Campus} besteht aus \glspl{Raum}n, die wiederum aus \glspl{Feld}n bestehen. Ein \gls{Feld} kann eine 
feste Anzahl an \glspl{GameObject} aufnehmen. Im \gls{Spiel} tauchen verschiedene \glspl{GameObject} auf, z.B. 
\glspl{Charakter}, \glspl{Objekt} oder \glspl{Gegenstand}. Sowohl der \gls{Spieler} als auch \gls{npcs} sind
\glspl{Charakter}. 
Das Spiel ist gewonnen, sobald die Gewinnbedingung erfüllt wurde. Die Gewinnbedingung besteht darin einen 
bestimmten \gls{Schein} zu erlangen. Ein \gls{Spiel} kann neu gestartet, abgespeichert oder geladen werden.
Das Ziel des \gls{Spiel}s könnte später auch durch einen Abschluss 
repräsentiert werden. Der Abschluss würde sich beispielsweise immer aus sechs Semestern zusammensetzen
und diese würden wiederum aus mehreren \glspl{Schein}n bestehen.

\section{Objekte und Gegenstände}
Grundsätzlich gibt es \glspl{Objekt} und \glspl{Gegenstand}. Objekte sind stationäre Gebilde auf einem Feld,
die der Spieler nicht mitnehmen kann. \glspl{Gegenstand} dagegen sind portabel, der Spieler kann sie in sein
Inventar packen. Auch NPCs können Gegenstände mit sich tragen. Häufig auftauchende Gegenstände sind zum 
Beispiel \glspl{Item}. Der Spieler soll mit jeder Form von Gegenstand und den meisten Objekte interagieren
können. Die Reaktion der Gegenstände auf bestimmte Interaktionen kann intern durch einfaches Nachschauen in
einer Tabelle geschehen. Jeder Gegenstand kann dann eine oder mehrere besondere Formen und eine standardmäßige
Form von Interaktion haben um alle Kombinationen möglichst einfach abdecken zu können.
In folgendem UML-Diagramm wird die Domainstruktur bildlich dargestellt. Die grünen Klassen zeigen mögliche 
spätere Erweiterungen an. Sie werden hier nur zur Orientierung abgebildet und sind nicht Teil der 
Spezifikation.
\newpage
\begin{figure}[htb]
	\begin{center}
		%trim option's parameter order: left bottom right top
		\includegraphics[trim = 5mm 5mm 20mm 40mm, clip, width=14cm]{kapitel/daten/domain.pdf}
	\end{center}
	\caption{UML-Diagramm der Domain}
	\label{fig:domain_uml}
\end{figure}

\newpage
\section{Inventar}
Das Inventar bietet dem \gls{Spieler} die Möglichkeit \glspl{Item} in seinen Besitz zu bringen. Es beinhaltet
eine feste Anzahl an Slots. Jeder Slot kann mit nur einem \gls{Item} belegt werden. Die Anzahl der Items im
Besitz des Spielers ist also durch die Kapazität des Inventars (plus eins im Handslot) direkt beschränkt.
Der Spieler kann Items aus dem Inventar ablegen (droppen) oder in die Hand nehmen (benutzen). Auch NPCs können
ein Inventar haben und Gegenstände mit sich herumtragen. Der Spieler kann Items also nicht nur durch das
Aufnehmen von Feldern, sondern auch durch die Interaktion (z.B. den Dialog) mit NPCs erlangen.

\section{Felder, Räume und Campi}
Felder sind die kleinste räumliche (ebene) Einheit im Spiel. Felder haben eine feste Anzahl Slots (eine 
Quadratzahl, in unserem Spiel neun) und können maximal Slots + 1 viele Elemente aufnehmen (der Spieler muss 
immer auf ein Feld
navigieren können, auch wenn alle Slots belegt sind). Jedes Feld hat kein bis maximal acht
angrenzende Felder. Felder haben Koordinaten. Der Übergang zwischen zwei Feldern muss nicht immer möglich
sein, so können diese durch eine Wand getrennt sein, was dem Spieler ersichtlich sein sollte.
Der Spieler steht immer genau auf einem Feld und kann alle unmittelbar angrenzenden Felder in Blickrichtung
sehen (aber 
nicht notwendigerweise erreichen oder betreten). Auf diesem Feld hat er eine von genau vier 
Blickrichtungen. Diese werden mit N, S, O, W bezeichnet. Diagonal angrenzende Felder kann er gegebenenfalls
sehen, jedoch nicht direkt betreten oder darauf zugreifen.
Räume bestehen aus zusammenhängenden Feldern, die einander zugänglich und vom Rest des \gls{Campus} durch
Wände abgetrennt sind. Räume sind mit einander durch Eingänge verbunden. Eingänge können verschlossen oder
offen sein. Eingänge können auch mit Items (z.B. Schlüsseln) kombiniert werden. Alle Räume zusammen bilden
einen Campus. In der vorgelegten Version spezifizieren wir genau einen Campus pro Spiel.

\section{NPCs und Dialoge}
NPCs sind Nichtspielercharaktere, die auf einem Feld stehen. Der Spieler kann mit ihnen in Form von Dialogen
oder durch Items interagieren. Dialoge haben keine oder mehrere Erwiderungen, die der Spieler dem NPC geben
kann. Eine Erwiderung kann einen Folgedialog haben. Dialoge werden also in Form eines Baums repräsentiert,
in der es einen Einstiegsdialog (Vaterknoten) mit Erwiderungen (den Pfaden) und Folgedialogen (Kindknoten)
geben kann. Manche Dialoge können Gegenstände enthalten, der Spieler erhält also den Gegenstand, vorausgesetzt
er wählt die richtige Abfolge von Erwiderungen. NPCs sollen stationär sein, in späteren Erweiterungen könnten
sich NPCs aber auch auf festen Routen durch die Map von Feld zu Feld bewegen.

\section{Spieler und Scheine}
In der hier vorgelegten Version gibt es nur einen Schein, den es zu finden gilt. Sobald der Schein gefunden
wurde, ist das Spiel automatisch beendet. In zukünftigen Erweiterungen könnten \glspl{Schein} dem 
\gls{Spieler} CPs einbringen. Das Finden der \glspl{Schein} könnte dann an eine Reihe von inhaltlichen
Verpflichtungen geknüpft sein (z.B. "finde \gls{Item} A, damit dir NPC B Schein C überreicht"). Das Ganze
könnte dann in Form von Semestern und Abschlüssen weiter ausgebaut werden, wie schon zuvor angedeutet.


\svnid{$Id: ui.tex 248 2012-06-17 17:10:59Z dgens001 $}
\pagebreak[4]
\chapter{Benutzungsschnittstellen}

\textit{Hier findet man die Dialogspezifikation, GUI-Skizzen, die Dialogabläufe und die für das Spiel wichtige
Formatspezifikation der Campusdatei.}

\section{Ingame GUI}
Ingame gibt es vier sog. \glspl{Aktivitaetsbereich}. Im \gls{Sichtbereich} (gelb) wird dem \gls{Spieler} seine 
aktuelle Position und die angrenzenden Felder angezeigt. Zudem werden alle 
\glspl{Item} und \gls{npcs} auf dem aktuellen Feld dargestellt. Nicht begehbare Felder sind durch Wände oder
verschlossene Türen abgetrennt. Offene Türen stellen betretbare Felder dar.
\\
\begin{figure}[h]
	\begin{center}
		\includegraphics[width=13cm, height=10cm]{kapitel/ui/ig_kreuz_small_color.jpg}
	\end{center}
	\caption{Die vier Aktivitätsbereiche}
	\label{fig:ig_kreuz_color}
\end{figure}

Im \gls{Infobereich} (orange) hat der Spieler entweder eine Übersicht der einzelnen Inventarslots 
(Inventarsicht) oder in Interaktion textuelle Informationen (Infosicht). Der \gls{Handbereich} (grün) zeigt 
dem Spieler an, ob er ein Objekt in 
der Hand hält und falls ja, welches. Der \gls{Statusbereich} (blau) zeigt ein Porträt des Charakters.
\\
\begin{figure}[h]
	\begin{center}
		\includegraphics[width=13cm, height=10cm]{kapitel/ui/ig_leer_small.jpg}
	\end{center}
	\caption{Ein leerer Sichtbereich}
	\label{fig:ig_leer}
\end{figure}

Das Inventar in Abbildung \ref{fig:ig_leer} enthält einige Items, die der Nutzer bereits aufgesammelt hat.
Er hat aber kein Objekt in der Hand, daher die Darstellung der leeren Hand. Da er vor einem freien Bereich
steht, ist die Sichtweite am größten. Selbst in diesem Fall kann der Spieler aber nur drei Felder weit sehen
(das auf dem er sich momentan befindet eingeschlossen). Dahinter befinden sich eventuell weitere Felder, die
aber durch eine Art Nebel verhangen sind. Der Horizont ist also nicht sichtbar. Darüber befindet sich der
Himmel bzw. innerhalb von Gebäuden die Decke. 
\newpage
\begin{figure}[h]
	\begin{center}
		\includegraphics[width=13cm, height=10cm]{kapitel/ui/ig_npc_small.jpg}
	\end{center}
	\caption{Interaktion mit einem NPC}
	\label{fig:ig_npc}
\end{figure}
Abbildung \ref{fig:ig_npc} zeigt die Interaktion mit einem NPC. Der Infobereich zeigt jetzt textuelle
Informationen zur Interaktion an. Auch Dialoge werden in dieser Form dargestellt. Der Status- und Handbereich
bleiben in dieser Sicht erhalten, was praktisch ist um zum Beispiel eine Tür mit einem Schlüssel zu öffnen:
Man hat direkt im Blick, ob der Schlüssel schon in der Hand liegt, also verfügbar für diese Interaktion ist.
Wäre dies nicht der Fall, müsste dieser erst aus dem Inventar in die Hand genommen werden.
\newpage
\section{Menü GUI}
Die Menüs stellen die im Abschnitt 3 dargestellten Funktionen zur Verfügung. Wie dort angesprochen startet
der Spieler im Hauptmenü. Die Menüs sollen alle im gleichen Stil angelegt werden, daher werden hier nur zwei
Menüskizzen gezeigt. Die einzelnen Einträge können wie gesagt aus Abschnitt 3 übernommen werden.

\begin{figure}[htb]
	\begin{center}
		\includegraphics[width=9cm, height=6cm]{kapitel/ui/menu_main_small.jpg}
	\end{center}
	\caption{Das Hauptmenü}
	\label{fig:menu_main}
\end{figure}
Laden- und Speicherscreen wird es eine tabellarische Sicht mit den \glspl{Savegame} geben. Man hat jeweils
die Möglichkeit eines auszuwählen und zu laden bzw. ein neues anzulegen oder einzelne \glspl{Savegame} zu 
löschen.
\begin{figure}[htb]
	\begin{center}
		\includegraphics[width=9cm, height=6cm]{kapitel/ui/menu_load_small.jpg}
	\end{center}
	\caption{Das Lademenü}
	\label{fig:menu_load}
\end{figure}

\newpage
\section{Campusdatei \textit{(v0.1)}}
Die Campusdatei ist ein grundlegendes Konzept im Spiel, um den Campus konfigurierbar zu machen.
Dabei gibt es eine Standardvorgehensweise: Den Campus in Form von ASCII-Zeichen zu codieren
und dann die so entstandene Textdatei zu parsen. Da das Erstellen eines neuen Campus von Hand
so aber sehr umständlich werden kann (und Probleme mit der Positionierung auftreten könnten),
sollte der Campus als Bitmapdatei codiert werden. Letzten Endes ist auch eine Bitmap eine Textdatei
in der die Pixel Felder im Spiel darstellen und kann geparsed werden. Entscheidender Vorteil ist die 
bessere Les- und somit Benutzbarkeit.

\begin{figure}[h]
	\begin{center}
		\includegraphics[width=7cm, height=7cm]{kapitel/ui/empty.png}
	\end{center}
	\caption{Ein leerer Campus}
	\label{fig:empty_campus}
\end{figure}

Das 29x29 Pixel große Minimalbeispiel einer solchen Datei wurde mit KolourPaint erstellt. Das Format der 
Datei ist PPM (Portable Pixel Map) in der Version P6 
(Spezifikation: \url{http://netpbm.sourceforge.net/doc/ppm.html}).
Die 3x3 großen Pixelfelder in Grau (\textit{0xC0C0C0}) stellen Felder auf dem Campus dar.
Die ein Pixel breiten Zwischenräume bieten Platz für Wände (\textit{0x000000}) und Türen (grün 
\textit{0x00FF00} oder rot \textit{0xFF0000}). Die Starposition des Spielers
wird durch einen blaues (\textit{0x0000FF}) Pixel innerhalb eines Feldes verdeutlicht. 

\begin{figure}[h]
	\begin{center}
		\includegraphics[width=7cm, height=7cm]{kapitel/ui/room_door_fancy.png}
	\end{center}
	\caption{Zwei Räume mit Tür}
	\label{fig:room_door_fancy}
\end{figure}

Es soll nicht möglich sein, Wände quer über ein Feld zu zeichnen. Allein die Zwischenräume bieten Platz für
Wanddefinitionen. Sollen Felder untereinander frei zugänglich sein, so bleibt der Zwischenraum weiss 
(\textit{0xFFFFFF}). Sollen Felder zwar durch eine Wand getrennt, aber mittels einer Tür verbunden sein,
so wird dies durch ein grünes (\textit{0x00FF00}) Pixel im Fall einer offenen, oder ein rotes 
(\textit{0xFF0000}) Pixel im Fall einer verschlossenen Tür dargestellt. Türen und Wände sind die einzigen
Gegenstände die im Zwischenraum positioniert werden dürfen. Türen belegen zusätzlich auf den Feldern die sie
verbinden einen Slot (in der gleichen Farbe die die jeweilige Tür markiert). Zu beachten ist auch, dass der
auf diese Art und Weise gezeichnete Campus immer einen ein Pixel großen Rand hat (\textit{0x000000} oder
\textit{0xFFFFFF}) und die Felder nur immer jeweils um drei Pixel gegeinander verschoben werden können, die
Feldkanten sind also immer bündig zu einander angelegt. Das gleiche gilt für Wände. Ein Feld ist immer 3x3
Pixel groß, ein Wandstück immer 3x1 Pixel lang. Es können allerdings mehrere dieser Wandteile aneinander
gefügt werden, wie in Abbildung \ref{fig:room_door_fancy} zu sehen, falls größere Bereiche zum Beispiel nicht
begehbar sein sollen.

Die Mapdatei muss eine Versionsnummer haben. Da die Versionsnummer für unser Spiel im PPM Format nicht 
vorgesehen ist, muss jede Mapdatei als Dateiendung "ctfmap\textit{xx}" haben. Die Nummer \textit{xx} ist die 
zugehörige Versionsnummer des Spiels (ohne trennenden Punkt) und die Campusdatei ist genau für diese Version
zulässig. Je nach Version können weitere Objekte auf
den Feldern der Map positioniert werden, Grundlage ist aber immer die Version 0.1. Die angegebenen Pixelgrößen
gehören zum definierten Campusformat dazu und müssen immer berücksichtigt werden. Dateien, die nicht genau das
angegebene Format haben, können nicht eingelesen werden.


% ---------------------------------- Anhänge -----------------------------------
% In diesem Teil werden alle Anhänge eingefügt, die auch als ganz normale Kapi-
% tel abgelegt werden.
%\appendix
%\include{kapitel/anhang}

% Ausgabe des Glossars (oder der Glossare, wenn mehrere definiert sind)
%\glsaddall
\printglossary

\end{document}
